\section{Recursos do Template}

% ----------------------------------------------------------------
\subsection{Código-Fonte com \texttt{minted}}

Para blocos com legenda e label, use \verb|listing| (float) + \verb|minted|:

\begin{listing}[H]
  \caption{Contador de 4 bits em VHDL}
  \label{cod:contador}
  \begin{minted}{vhdl}
library IEEE;
use IEEE.STD_LOGIC_1164.ALL;
use IEEE.NUMERIC_STD.ALL;

entity contador is
  port (
    clk   : in  std_logic;
    reset : in  std_logic;
    q     : out std_logic_vector(3 downto 0)
  );
end entity;

architecture rtl of contador is
  signal count : unsigned(3 downto 0) := (others => '0');
begin
  process(clk)
  begin
    if rising_edge(clk) then
      if reset = '1' then
        count <= (others => '0');
      else
        count <= count + 1;
      end if;
    end if;
  end process;
  q <= std_logic_vector(count);
end architecture;
  \end{minted}
\end{listing}

\begin{listing}[H]
  \caption{Script de análise de simulação}
  \label{cod:python}
  \begin{minted}{python}
import numpy as np
import matplotlib.pyplot as plt

# Lê saída do GTKWave exportada como CSV
data = np.loadtxt("sim.csv", delimiter=",", skiprows=1)
t, q = data[:, 0], data[:, 2]

plt.plot(t * 1e9, q, drawstyle="steps-post")
plt.xlabel("Tempo (ns)")
plt.ylabel("Saída q")
plt.savefig("forma_de_onda.pdf")
  \end{minted}
\end{listing}

Para trechos inline use \verb|\mintinline|,
por exemplo: \mintinline{vhdl}{if rising_edge(clk) then}.

Referências: \cref{cod:contador} e \cref{cod:python}.

% ----------------------------------------------------------------
\subsection{Unidades Físicas com \texttt{siunitx}}

Use \verb|\qty{valor}{unidade}| para quantidades e \verb|\num{valor}| para números.
A vírgula decimal e o separador de milhar seguem o padrão brasileiro automaticamente.

Exemplos de quantidades comuns em Engenharia Elétrica:

\begin{itemize}
  \item Frequência de clock: \qty{50}{\mega\hertz}
  \item Tensão de alimentação: \qty{3.3}{\volt}
  \item Corrente de fuga: \qty{1.5}{\micro\ampere}
  \item Período de amostragem: \qty{20}{\nano\second}
  \item Resistência: \qty{10}{\kilo\ohm}
  \item Capacitância: \qty{100}{\pico\farad}
\end{itemize}

Número com separador de milhar: \num{1234567}.
Notação científica: \num{1.5e-9}.

\tabeladuas
  {Parâmetros de temporização}
  {tab:timing}
  {Parâmetro}{Valor}
  {
    Período de clock  & \qty{20}{\nano\second} \\
    Setup time        & \qty{1.5}{\nano\second} \\
    Hold time         & \qty{0.25}{\nano\second} \\
    Tensão VCC        & \qty{3.3}{\volt} \\
    Corrente máxima   & \qty{200}{\milli\ampere} \\
  }

% ----------------------------------------------------------------
\subsection{Subfiguras com \texttt{subcaption}}

Use o ambiente \verb|subfigure| para colocar figuras lado a lado com legendas
e rótulos independentes. Substitua os retângulos abaixo por \verb|\includegraphics|.

\begin{figure}[H]
  \centering
  \begin{subfigure}[b]{0.45\textwidth}
    \centering
    \begin{tikzpicture}
      \draw[fill=azulpoli!15, rounded corners=4pt] (0,0) rectangle (4,2.5);
      \node at (2,1.25) {\small Imagem A};
    \end{tikzpicture}
    \caption{Primeira subfigura.}
    \label{fig:sub-a}
  \end{subfigure}
  \hfill
  \begin{subfigure}[b]{0.45\textwidth}
    \centering
    \begin{tikzpicture}
      \draw[fill=azulpoli!35, rounded corners=4pt] (0,0) rectangle (4,2.5);
      \node at (2,1.25) {\small Imagem B};
    \end{tikzpicture}
    \caption{Segunda subfigura.}
    \label{fig:sub-b}
  \end{subfigure}
  \caption{Exemplo de subfiguras com \texttt{subcaption}.}
  \label{fig:subfiguras}
\end{figure}

% ----------------------------------------------------------------
\subsection{Referências Automáticas com \texttt{cleveref}}

Use \verb|\cref{label}| no lugar de escrever \verb|Figura~\ref{label}| na mão.
O tipo da referência é inserido automaticamente em português:

\begin{itemize}
  \item \verb|\cref{fig:subfiguras}|          $\to$ \cref{fig:subfiguras}
  \item \verb|\cref{fig:sub-a}|               $\to$ \cref{fig:sub-a}
  \item \verb|\cref{fig:sub-a,fig:sub-b}|     $\to$ \cref{fig:sub-a,fig:sub-b}
  \item \verb|\cref{tab:timing}|              $\to$ \cref{tab:timing}
  \item \verb|\cref{eq:integral}|             $\to$ \cref{eq:integral}
\end{itemize}

\verb|\Cref| (maiúsculo) é usado no início de frase:
\Cref{fig:subfiguras} mostra um exemplo de subfiguras lado a lado.
